%--------------------------------------------------------------------------------------
% AI Nyilatkozat
%--------------------------------------------------------------------------------------
%----------------------------------------------------------------------------
\section{Nyilatkozat generatív mesterséges intelligencia alkalmazásáról}
%----------------------------------------------------------------------------

\noindent $\Box$ \textbf{Nem használtam} semmilyen generatív MI segédeszközt.

\noindent $\boxtimes$ \textbf{Használtam} generatív MI segédeszközt. Az MI-vel generált tartalmakat ellenőriztem, a generált kimenetek valóságtartalmáról meggyőződtem, az alábbi táblázatban megfelelően jelöltem minden használatot.

\small{\noindent
	\begin{tabular}{|>{\raggedright\arraybackslash}p{38mm}|>{\raggedright\arraybackslash}p{30mm}|>{\raggedright\arraybackslash}p{36mm}|>{\raggedright\arraybackslash}p{30mm}|}\hline
		\textbf{Felhasználási módok}	& \textbf{Generatív MI eszköz(ök) neve} 	& \textbf{Érintett részek} (fejezet, oldalszám, hivatkozás) 	& \textbf{Használat becsült aránya} (felhasználási módonként) \\ \hline \hline
		A \LaTeX-forrás szerkesztése, formázása & Claude (Anthropic), Claude Code & teljes dolgozat & \textbf{100}\% \\ \hline
		\textbf{Prompt lényegi része}  &\multicolumn{3}{p{100mm}|}{A dolgozat \LaTeX-forrásának szerkesztése: ábrák elrendezése és méretezése, tördelés, hivatkozások és képaláírások javítása.} \\ \hline \hline
		Adatvizualizáció, grafikonok generálása adatpontok alapján & Claude (Anthropic), Claude Code & \ref{fig:amazonHeatmap}. és \ref{fig:versions}.~ábra & 100\% \\ \hline
		\textbf{Prompt lényegi része}  &\multicolumn{3}{p{100mm}|}{Grafikonok és hőtérképek készítése meglévő mérési adatokból: a verziók közötti mérések grafikonja és az Amazon hőtérképei.} \\ \hline
	\end{tabular}
} % end of \small
